% !TeX encoding = UTF-8
% !TeX root = thesis.tex

\section{Grundlagen und verwandte Arbeiten}

% ---------------------------------------------------------------------------
\subsection{Cloud-Anomalieerkennung und Alert Fatigue}
% ---------------------------------------------------------------------------

Moderne Cloud-Umgebungen werden kontinuierlich über systemseitige Metriken wie
CPU-Auslastung, Speicherverbrauch und Netzwerkdurchsatz überwacht. Anomalie-Detektoren
werten diese Zeitreihendaten automatisiert aus und generieren Alerts bei auffälligen
Abweichungen. In großen, dynamischen Infrastrukturen erzeugen solche Systeme jedoch
regelmäßig eine hohe Zahl von Fehlalarmen, sogenannte False Positives, die
Sicherheitsanalysten belasten und deren Reaktionsfähigkeit auf echte Vorfälle
verringern. Dieser Effekt ist in der Literatur als \emph{Alert Fatigue} bekannt
und stellt ein anerkanntes operatives Problem in Security Operations Centers (SOCs) dar~\cite{tariq2025alert}.
Alahmadi~et~al.~\cite{alahmadi202299} belegen in einer qualitativen Studie, dass
in real betriebenen SOCs ein erheblicher Teil der ausgelösten Sicherheitsalarme auf
False Positives entfällt, was die Analysekapazität für tatsächliche Vorfälle
systematisch einschränkt.

Die False-Positive-Rate (FPR) ist dabei die zentrale operative Zielgröße: Sie
beschreibt den Anteil tatsächlich negativer Instanzen, der fälschlicherweise als
Alarm eingestuft wird. Eine hohe FPR verursacht nicht nur unnötigen Analyseaufwand,
sondern kann im Extremfall dazu führen, dass echte Angriffe im Rauschen untergehen.
Die Reduktion der FPR bei gleichzeitiger Aufrechterhaltung einer hohen
Erkennungsrate (Recall) bildet daher die Forschungsfrage dieser Arbeit.

Als Evaluierungsgrundlage dient CloudAnoBench~\cite{zou2025cloudanobench}, ein
synthetischer Benchmark, der Systemmetriken und Log-Dateien für Szenarien mit
malignem, anomalem und normalem Systemverhalten bereitstellt. CloudAnoBench ist
ein neuerer Benchmark und nicht als langjährig etablierter Goldstandard der
Community zu verstehen; die in dieser Arbeit erzielten Ergebnisse sind
datensatzspezifisch und besitzen keine unmittelbare Allgemeingültigkeit.

% ---------------------------------------------------------------------------
\subsection{Metrikbasierte und logbasierte Anomalieerkennung}
% ---------------------------------------------------------------------------

Anomalieerkennung in Cloud-Systemen kann grundsätzlich auf zwei komplementären
Informationsquellen operieren: Systemmetriken und Log-Dateien.

\textbf{Systemmetriken} sind numerische Zeitreihendaten, die den Ressourcenzustand
eines Systems zu diskreten Zeitpunkten quantifizieren, etwa Prozessorauslastung,
Speicherverbrauch oder Netzwerkdurchsatz. Sie sind gut strukturiert, hochfrequent
verfügbar und lassen sich effizient mit tabellarischen Lernverfahren verarbeiten.
Ihre Schwäche besteht darin, dass unterschiedliche Ereignistypen, etwa legitime
Lastspitzen und Angriffsmuster, identische Metrikverläufe erzeugen können. Metriken
allein erlauben in solchen Fällen keine zuverlässige Unterscheidung.

\textbf{Log-Dateien} hingegen enthalten textuelle Beschreibungen von Systemereignissen,
Prozessaufrufen und Fehlerzuständen. Sie liefern einen prozessualen Kontext, der über
die bloße Ressourcensicht hinausgeht. He~et~al.~\cite{he2021survey} zeigen in einer
umfassenden Übersicht, dass automatisierte Log-Analyse ein etabliertes Werkzeug im
Reliability Engineering ist: Textmuster in Log-Dateien können charakteristische
Signaturen für unterschiedliche Ereigniskategorien tragen, die in den Metriken nicht
sichtbar sind.

Die vorliegende Arbeit verfolgt einen zweistufigen Ansatz: Ein metrischer Detektor
identifiziert zunächst potenzielle Auffälligkeiten auf Basis von Ressourcensignalen;
eine log-basierte Komponente verifiziert anschließend auf Case-Ebene, ob der
Log-Kontext den Alert stützt. Beide Komponenten sind methodisch aufeinander abgestimmt
und werden in Abschnitt~\ref{sec:data_preparation} im Detail beschrieben.

% ---------------------------------------------------------------------------
\subsection{Unbalancierte Klassifikation und Precision-Recall-Metriken}
% ---------------------------------------------------------------------------

Anomalie-Erkennungsprobleme sind typischerweise stark unbalanciert: Positive
Instanzen (Angriffe, Fehler) sind weit seltener als negative (Normalbetrieb).
He und Garcia~\cite{hegarcia2009learning} zeigen, dass in solchen Szenarien
standardmäßige Klassifikationsmetriken wie Accuracy irreführend sind, da sie
von der Mehrheitsklasse dominiert werden.

Die für diese Arbeit relevanten Metriken sind:

\begin{description}
  \item[Precision] $P = \frac{TP}{TP + FP}$: Anteil der korrekt als positiv
    klassifizierten Instanzen an allen als positiv gekennzeichneten.
  \item[Recall] $R = \frac{TP}{TP + FN}$: Anteil der tatsächlich positiven
    Instanzen, die erkannt werden. In dieser Arbeit dient $R \geq 0{,}95$ als
    Zielkriterium für die Threshold-Wahl auf dem Validierungsset; ob dieser Wert
    im unabhängigen Testsplit erreicht wird, ist eine empirische Frage.
  \item[False-Positive-Rate] $\text{FPR} = \frac{FP}{FP + TN}$: Anteil der
    negativen Instanzen, die fälschlicherweise als positiv eingestuft werden.
    Dies ist die zentrale operative Zielgröße.
  \item[Average Precision (AP)] Fläche unter der Precision-Recall-Kurve,
    berechnet als gewichteter Mittelwert der Precision-Werte über alle
    Recall-Stufen. AP fasst die Klassifikatorleistung über alle möglichen
    Thresholds zusammen.
\end{description}

Davis und Goadrich~\cite{davis2006relationship} zeigen, dass die PR-Kurve
gegenüber der ROC-Kurve bei unbalancierten Problemen informativer ist: Während
ROC-AUC die Leistung im negativen Regime mitunter überschätzt, macht die PR-Kurve
Schwächen bei der Behandlung der Minderheitsklasse direkt sichtbar. Saito und
Rehmsmeier~\cite{saito2015precision} bestätigen dies und empfehlen PR-Kurven
explizit für stark unbalancierte Klassifikationsprobleme. Diese Befunde motivieren
in dieser Arbeit die Verwendung von AUCPR als primäre Optimierungsmetrik sowie
die Fokussierung auf FPR und AP in der Evaluation.

% ---------------------------------------------------------------------------
\subsection{XGBoost für tabellarische Metrikdaten}
% ---------------------------------------------------------------------------

Als metrischer Basis-Detektor wird in dieser Arbeit XGBoost~\cite{chen2016xgboost}
eingesetzt. XGBoost implementiert Gradient Tree Boosting und approximiert den
Split-Finding-Algorithmus durch eine histogrammbasierte Methode
(\texttt{tree\_method='hist'}), die den Trainingsaufwand gegenüber dem exakten
Verfahren deutlich reduziert. Für tabellarisch strukturierte Daten mit numerischen
Features und heterogenen Skalen stellt XGBoost ein in der Literatur breit erprobtes
Verfahren dar.

Relevante Eigenschaften für diese Arbeit sind der Parameter
\texttt{scale\_pos\_weight} zur Gewichtung der Minderheitsklasse im unbalancierten
Setting sowie die Möglichkeit, AUCPR direkt als Optimierungsmetrik zu verwenden.

XGBoost wird hier nicht als grundsätzlich überlegenes Verfahren positioniert, sondern
als solide, reproduzierbare Baseline für tabellarische Klassifikation. Es ist nicht
Ziel dieser Arbeit, das metrische Modell zu optimieren; vielmehr soll seine
Leistungsgrenze die Notwendigkeit einer ergänzenden Log-Komponente motivieren.

% ---------------------------------------------------------------------------
\subsection{TF-IDF und logistische Regression für Logdaten}
% ---------------------------------------------------------------------------

Die log-basierte Komponente dieser Arbeit nutzt eine Kombination aus
\emph{TF-IDF}-Vektori\-sie\-rung und logistischer Regression. TF-IDF (\emph{Term
Frequency~-- Inverse Document Frequency}) gewichtet Terme proportional zu ihrer
Häufigkeit im Dokument, normiert jedoch durch ihre Häufigkeit über alle Dokumente.
Häufige, wenig diskriminierende Terme, in Log-Dateien etwa routinemäßige
Statusmeldungen, erhalten dadurch geringere Gewichte; seltene, charakteristische
Terme werden hervorgehoben.

Auf der TF-IDF-Matrix wird eine logistische Regression trainiert. Lineare Modelle
auf sparse hochdimensionalen Textrepräsentationen sind für diese Problemklasse
eine etablierte Wahl; umgesetzt wird der Klassifikator in dieser Arbeit mit
scikit-learn~\cite{pedregosa2011scikit}. Die logistische Regression
liefert Scores, die in der nachgelagerten Fusionsstufe
direkt weiterverwendet werden. Eine explizite Wahrscheinlichkeitskalibrierung
(z.\,B. Platt Scaling) wurde nicht durchgeführt; die Ausgaben werden daher als
Modellscores interpretiert.

Die Entscheidung für TF-IDF und logistische Regression folgt dem Prinzip einer
interpretierbaren und effizient trainierbaren Baseline. Fortgeschrittenere Ansätze zur Log-Analyse, etwa struktur-bewusste
Log-Parser, vortrainierte Sprachmodelle oder Embedding-Methoden, wurden in dieser
Arbeit nicht untersucht. Sie stellen eine naheliegende Erweiterung dar, sind aber
nicht Gegenstand der vorliegenden Evaluation.

% ---------------------------------------------------------------------------
\subsection{Data Leakage und gruppenbasierte Validierung}
% ---------------------------------------------------------------------------

Eine korrekte Evaluation maschineller Lernverfahren setzt voraus, dass keine
Information aus dem Testset in den Trainingsprozess einfließt. Kaufman~et~al.~\cite{kaufman2012leakage}
klassifizieren verschiedene Formen von Datenleck (\emph{Data Leakage}) und zeigen,
dass selbst scheinbar harmlose Preprocessing-Schritte, etwa eine globale
Skalierung oder Imputation, die gemessene Testleistung künstlich erhöhen können,
wenn sie auf dem Gesamtdatensatz statt ausschließlich auf dem Trainingsanteil
durchgeführt werden.

Für diese Arbeit ist eine spezifische Form von Leakage besonders relevant:
\emph{Schema-Leakage} durch strukturelle Fehlwerte. Werden verschiedene Klassen
mit systematisch unterschiedlichen Feature-Schemata erfasst, entstehen
typspezifische \texttt{NaN}-Muster, die ein Modell implizit zur Bestimmung der
Klassenzugehörigkeit nutzen kann, noch bevor inhaltlich relevante Signale
verarbeitet werden. Kaufman~et~al.~\cite{kaufman2012leakage} klassifizieren dies
als Target Leakage, da die Feature-Struktur die Zielvariable enkodiert und eine
valide Evaluation verhindert. Die methodische Konsequenz ist die Beschränkung der
Modellierung auf Features, die in allen Klassen vollständig verfügbar sind.

Ein zweites Leakage-Risiko ergibt sich aus der Datensatzstruktur: Beobachtungen
innerhalb desselben Szenarios sind strukturell korreliert. Eine zufällige
Aufteilung in Trainings- und Testmenge ohne Berücksichtigung dieser Gruppen
würde Szenarien auf beiden Seiten des Splits erlauben und die Trennleistung
systematisch überschätzen. Roberts~et~al.~\cite{roberts2017cross} zeigen, dass
bei Daten mit gruppenweise korrelierten Beobachtungen eine gruppenweise Partitionierung
notwendig ist. In dieser Arbeit wird \texttt{GroupShuffleSplit} verwendet, das
sicherstellt, dass kein Szenario gleichzeitig in Trainings- und Testmenge vertreten
ist.

% ---------------------------------------------------------------------------
\subsection{Klassifikatorfusion und Case-Level-Evaluation}
% ---------------------------------------------------------------------------

Die Kombination mehrerer Klassifikatoren zu einem gemeinsamen Entscheidungssystem
ist ein etabliertes Forschungsfeld. Kittler~et~al.~\cite{kittler1998combining}
systematisieren verschiedene Fusionsregeln, darunter Produkt-Regel, Summen-Regel
und Mehrheitsentscheid, und zeigen, unter welchen Voraussetzungen sie
asymptotisch äquivalent sind.

In dieser Arbeit wird die Fusionsentscheidung direkt auf Case-Ebene getroffen.
Beide Komponenten liefern pro Case~$c$ einen Score:
$\hat{p}_{\text{metric}}(c)$ aus dem metrischen XGBoost-Klassifikator, der auf
25~Case-Level-Aggregaten der fünf universellen Systemmetriken trainiert wird
(je Metrik: Mittelwert, Maximum, Standardabweichung, 95.~Perzentil und linearer
Trend), sowie $\hat{p}_{\text{log}}(c)$ aus der log-basierten logistischen
Regression auf dem normalisierten TF-IDF-Raum. Beide Komponenten arbeiten direkt
auf Case-Ebene; ein zwischengeschalteter Row-Level-Aggregationsschritt entfällt.

Zwei Fusionsstrategien werden untersucht. Das \textbf{AND-Gate} stellt eine
konservative Strategie dar: Ein Alert wird nur ausgegeben, wenn \emph{beide}
Komponenten positiv entscheiden:
\[
  \hat{y}_{\text{AND}}(c) = \mathbf{1}\!\bigl[
    \hat{p}_{\text{metric}}(c) \geq \tau_{\text{metric}}
    \;\wedge\;
    \hat{p}_{\text{log}}(c) \geq \tau_{\text{log}}
  \bigr]
\]
Die \textbf{Soft-Fusion} kombiniert beide Scores als Produkt:
\[
  s_{\text{fusion}}(c) = \hat{p}_{\text{metric}}(c) \times \hat{p}_{\text{log}}(c)
\]
Da beide Faktoren im Intervall $[0,1]$ liegen, gilt dies auch für
$s_{\text{fusion}}(c)$. Kittler~et~al.~\cite{kittler1998combining} zeigen, dass
die Produktregel unter der Annahme unabhängiger Klassifikatoren asymptotisch
optimal ist; wie gut diese Annahme auf den vorliegenden Daten zutrifft, ist eine
empirische Frage. Die Thresholds $\tau_{\text{metric}}$ und $\tau_{\text{log}}$
werden ausschließlich auf dem Validierungsset bestimmt; das Testset wird nur zur
finalen Evaluation verwendet. Ob die AND-Gate-Strategie gegenüber einer
Einzelkomponente einen Mehrwert liefert, wird in Abschnitt~\ref{sec:ablation}
untersucht.
